%% ============================================================
%%  TEMPLATE PPGCC UFERSA/UERN  –  dissertacao.tex
%%  Versão 2.0 (2025)  |  Normas ABNT NBR 14724:2011
%% ============================================================
%%
%%  COMO USAR:
%%  1. Preencha os metadados na seção "CONFIGURAÇÕES DO TRABALHO"
%%  2. Edite os arquivos em  capitulos/   (textuais)
%%  3. Edite os arquivos em  pre-textuais/
%%  4. Edite pos-textuais/referencias.bib
%%  5. Compile com:  pdflatex → bibtex → pdflatex → pdflatex
%%     (ou use o Makefile: make)
%%
%% ============================================================

\documentclass[
  a4paper,           % papel A4
  12pt,              % fonte 12pt (ABNT)
  chapter=TITLE,     % capítulos em CAIXA ALTA
  section=Title,     % seções com Inicial Maiúscula
  subsection=Title,
  oneside,           % impressão em uma face
  english,
  brazil             % idioma principal (deve ser o último)
]{abntex2}

%% ── CARREGA O ESTILO PPGCC ───────────────────────────────────
\usepackage{lib/ppgcc}

%% ============================================================
%%  CONFIGURAÇÕES DO TRABALHO  ← edite aqui
%% ============================================================

%% Tipo: dissertacao | tese
\ppgccTipoTrabalho{dissertacao}

%% Coloque 'sim' durante a qualificação, 'nao' para versão final
\ppgccEhQualificacao{nao}

%% ── Informações institucionais ───────────────────────────────
\ppgccIES{Universidade Federal Rural do Semi-Árido}
\ppgccIESDois{Universidade Estadual do Rio Grande do Norte}
\ppgccPrograma{Programa de Pós-Graduação em Ciência da Computação}
\ppgccArea{Ciência da Computação}

%% ── Dados do trabalho ────────────────────────────────────────
\titulo{Título da Dissertação: Subtítulo se Houver}
\autor{Nome Completo do Autor}
\ppgccLocal{Mossoró}
\ppgccAno{2025}
\ppgccDataAprovacao{01 de janeiro de 2025}

%% ── Orientação ───────────────────────────────────────────────
\ppgccOrientador{Prof. Dr. Nome do Orientador}
\ppgccOrientadorIES{UFERSA}
\ppgccOrientadorFem{nao}   % 'sim' se orientadora

%% Coorientador — deixe em branco ({}) para omitir
\ppgccCoorientador{Prof. Dr. Nome do Coorientador}
\ppgccCoorientadorIES{UERN}
\ppgccCoorientadorFem{nao}

%% ── Banca examinadora ────────────────────────────────────────
%% Membro 2 (obrigatório)
\ppgccBancaDois{Prof. Dr. Segundo Membro da Banca}
\ppgccBancaDoisIES{Universidade do Segundo Membro}

%% Membro 3 (obrigatório)
\ppgccBancaTres{Prof. Dr. Terceiro Membro da Banca}
\ppgccBancaTresIES{Universidade do Terceiro Membro}

%% Membro 4 (opcional — deixe {} para omitir)
\ppgccBancaQuatro{}
\ppgccBancaQuatroIES{}

%% ── Palavras-chave (para o resumo) ───────────────────────────
% Defina aqui e use \@palavraschave nos arquivos de resumo
\newcommand{\palavraschave}{Palavra-chave 1. Palavra-chave 2. Palavra-chave 3.}
\newcommand{\keywords}{Keyword 1. Keyword 2. Keyword 3.}

%% ============================================================
%%  INÍCIO DO DOCUMENTO
%% ============================================================
\begin{document}

%% ────────────────────────────────────────────────────────────
%%  ELEMENTOS PRÉ-TEXTUAIS
%% ────────────────────────────────────────────────────────────
\pretextual

%% Capa (obrigatório)
\imprimircapa

%% Folha de rosto (obrigatório)
\imprimirfolhaderosto{}

%% Ficha catalográfica — gerada pela biblioteca, inserir como PDF
%% Descomente quando tiver o arquivo:
% \begin{fichacatalografica}
%   \includepdf{pre-textuais/ficha-catalografica.pdf}
% \end{fichacatalografica}

%% Errata (opcional) — descomente se necessário
% %% pre-textuais/errata.tex
%% Elemento opcional. Folha solta após a folha de rosto.

\begin{errata}

\begin{tabular}{|p{2cm}|p{2cm}|p{5cm}|p{5cm}|}
\hline
\textbf{Folha} & \textbf{Linha} & \textbf{Onde se lê} & \textbf{Leia-se} \\
\hline
XX & XX & texto incorreto & texto correto \\
\hline
\end{tabular}

\end{errata}


%% Folha de aprovação (obrigatório após defesa)
\imprimirfolhadeaprovacao

%% Dedicatória (opcional)
%% pre-textuais/dedicatoria.tex
%% Elemento opcional. Recuo de 8 cm da esquerda para o centro.

\begin{dedicatoria}
   \vspace*{\fill}
   \hspace{.45\textwidth}
   \begin{minipage}{.5\textwidth}
      \textit{Dedique este trabalho a quem você quiser.
      Por exemplo, à minha família, pelo apoio incondicional
      durante toda essa jornada.}
   \end{minipage}
\end{dedicatoria}


%% Agradecimentos (opcional)
%% pre-textuais/agradecimentos.tex
%% Elemento opcional. Espaçamento 1,5. Sem recuo de parágrafo.

\begin{agradecimentos}

Agradeço a todos que contribuíram para a realização deste trabalho.

Em especial, agradeço ao meu orientador, Prof. Dr. Nome do Orientador,
pela dedicação, paciência e orientação ao longo de todo o processo.

Agradeço também ao Programa de Pós-Graduação em Ciência da Computação
da UFERSA/UERN e a todos os professores que contribuíram para a minha
formação.

Por fim, agradeço à minha família e amigos pelo apoio e compreensão
durante os momentos difíceis.

\end{agradecimentos}


%% Epígrafe (opcional)
%% pre-textuais/epigrafe.tex
%% Elemento opcional. Citação relacionada ao tema. Recuo de 8 cm.

\begin{epigrafe}
    \vspace*{\fill}
    \begin{flushright}
        \textit{``Escreva aqui uma citação relacionada ao tema
        da sua pesquisa.''}\\
        (Autor da Citação)
    \end{flushright}
\end{epigrafe}


%% Resumo em português (obrigatório)
%% pre-textuais/resumo.tex
%% Obrigatório. 150–500 palavras. Parágrafo único. Voz ativa, 3ª pessoa.
%% Estrutura: objetivo → método → resultados → conclusões.

\begin{resumo}

Este trabalho apresenta [objetivo geral do trabalho]. O contexto
que motiva a pesquisa é [contextualização do problema]. O problema
de pesquisa consiste em [descrição do problema]. Para isso,
[descreva o método utilizado: tipo de pesquisa, paradigma
metodológico, instrumentos]. Os principais resultados obtidos
demonstram que [síntese dos resultados]. Conclui-se que
[conclusão principal e contribuições].

\vspace{\onelineskip}

\noindent
\textbf{Palavras-chave:} \palavraschave

\end{resumo}


%% Resumo em inglês — Abstract (obrigatório para dissertação/tese)
%% pre-textuais/abstract.tex
%% Obrigatório para dissertação/tese. Mesmo padrão do resumo.

\begin{resumo}[Abstract]
  \begin{otherlanguage*}{english}

This work presents [general objective of the work]. The context
that motivates the research is [contextualization of the problem].
The research problem consists of [description of the problem].
To this end, [describe the method used]. The main results demonstrate
that [synthesis of results]. It is concluded that [main conclusion
and contributions].

\vspace{\onelineskip}

\noindent
\textbf{Keywords:} \keywords

  \end{otherlanguage*}
\end{resumo}


%% Lista de figuras (opcional; recomendado se ≥ 5 figuras)
\pdfbookmark[0]{\listfigurename}{lof}
\listoffigures*
\cleardoublepage

%% Lista de tabelas (opcional; recomendado se ≥ 5 tabelas)
\pdfbookmark[0]{\listtablename}{lot}
\listoftables*
\cleardoublepage

%% Lista de abreviaturas e siglas (opcional)
%% pre-textuais/abreviaturas.tex
%% Elemento opcional. Lista alfabética das siglas utilizadas no texto.

\begin{siglas}
  \item[ABNT]   Associação Brasileira de Normas Técnicas
  \item[API]    Application Programming Interface
  \item[LIMS]   Laboratory Information Management System
  \item[PPGCC]  Programa de Pós-Graduação em Ciência da Computação
  \item[RBAC]   Role-Based Access Control
  \item[SUS]    System Usability Scale
  \item[UFERSA] Universidade Federal Rural do Semi-Árido
  \item[UERN]   Universidade Estadual do Rio Grande do Norte
  %% Adicione mais siglas aqui em ordem alfabética
\end{siglas}


%% Sumário (obrigatório)
\pdfbookmark[0]{\contentsname}{toc}
\tableofcontents*
\cleardoublepage

%% ────────────────────────────────────────────────────────────
%%  ELEMENTOS TEXTUAIS
%% ────────────────────────────────────────────────────────────
\textual

%% capitulos/1-introducao.tex

\chapter{Introdução}
\label{cap:introducao}

Este capítulo apresenta a contextualização do problema de pesquisa,
a motivação, os objetivos, a metodologia adotada e a estrutura
desta dissertação.

% ── Dica: use \section, \subsection etc. para subdivisões

\section{Contextualização e Motivação}
\label{sec:contextualizacao}

Descreva aqui o contexto da área e o que motiva a pesquisa.
Cite autores relevantes \cite{exemplo2023}.

\section{Problema de Pesquisa}
\label{sec:problema}

Enuncie o problema de pesquisa de forma clara e objetiva.

\section{Objetivos}
\label{sec:objetivos}

\subsection{Objetivo Geral}

O objetivo geral desta dissertação é \ldots

\subsection{Objetivos Específicos}

Para atingir o objetivo geral, os seguintes objetivos específicos
foram definidos:

\begin{enumerate}
  \item Objetivo específico 1;
  \item Objetivo específico 2;
  \item Objetivo específico 3.
\end{enumerate}

\section{Metodologia}
\label{sec:metodologia-intro}

% Breve descrição da abordagem metodológica (detalhada no cap. 4)
Esta pesquisa adota o paradigma da Design Science Research
\cite{Hevner2004}, \ldots

\section{Contribuições}
\label{sec:contribuicoes}

As principais contribuições deste trabalho são:

\begin{itemize}
  \item Contribuição 1;
  \item Contribuição 2.
\end{itemize}

\section{Estrutura da Dissertação}
\label{sec:estrutura}

Este trabalho está organizado da seguinte forma:
o Capítulo~\ref{cap:fundamentacao} apresenta a fundamentação teórica;
% Adicione os demais capítulos aqui

%% capitulos/2-fundamentacao-teorica.tex

\chapter{Fundamentação Teórica}
\label{cap:fundamentacao}

Este capítulo apresenta os conceitos teóricos que fundamentam
esta dissertação.

\section{Tema Principal}
\label{sec:tema-principal}

Apresente os conceitos centrais da sua pesquisa, com embasamento
bibliográfico adequado \cite{exemplo2023}.

% Exemplo de citação direta curta (≤ 3 linhas, entre aspas):
Segundo \citeonline{exemplo2023}, ``texto citado diretamente entre
aspas, com até três linhas, inserido no corpo do parágrafo''.

% Exemplo de citação direta longa (> 3 linhas):
% Use o ambiente 'citacao' do abntex2:
\begin{citacao}
Texto com mais de três linhas que deve ser recuado 4 cm da margem
esquerda, com fonte tamanho 10, espaço simples, sem aspas.
O recuo e formatação são automáticos com este ambiente.
\cite[p. 15]{exemplo2023}
\end{citacao}

\section{Subtema 1}
\label{sec:subtema1}

Texto da seção \ldots

\subsection{Subseção 1.1}
\label{subsec:subsecao11}

Texto da subseção \ldots

% Exemplo de figura:
\begin{figure}[htb]
  \caption{Título descritivo da figura}
  \label{fig:exemplo}
  \begin{center}
    \includegraphics[width=0.6\textwidth]{figuras/exemplo}
  \end{center}
  \legend{Fonte: Autor (2025)}
\end{figure}

A Figura~\ref{fig:exemplo} ilustra \ldots

% Exemplo de tabela:
\begin{table}[htb]
  \caption{Título descritivo da tabela}
  \label{tab:exemplo}
  \begin{center}
    \begin{tabular}{lcc}
      \toprule
      \textbf{Coluna 1} & \textbf{Coluna 2} & \textbf{Coluna 3} \\
      \midrule
      Dado A & 10 & 20 \\
      Dado B & 15 & 30 \\
      \bottomrule
    \end{tabular}
  \end{center}
  \legend{Fonte: Elaborada pelo autor (2025)}
\end{table}

A Tabela~\ref{tab:exemplo} apresenta \ldots

%% capitulos/3-trabalhos-relacionados.tex

\chapter{Trabalhos Relacionados}
\label{cap:relacionados}

Este capítulo apresenta um mapeamento da literatura sobre os temas
centrais desta pesquisa, identificando lacunas que este trabalho
busca preencher.

\section{Critérios de Revisão}

Descreva como a revisão foi conduzida (strings de busca, bases
consultadas, critérios de inclusão/exclusão).

\section{Síntese dos Trabalhos}

% Organize por tema ou cronologicamente

\section{Comparativo e Lacunas}

% Tabela comparativa dos trabalhos relacionados

\section{Considerações do Capítulo}

Síntese das lacunas identificadas e como este trabalho as aborda.

%% capitulos/4-metodologia.tex

\chapter{Metodologia}
\label{cap:metodologia}

Descreva o paradigma metodológico adotado (ex.: Design Science
Research, pesquisa-ação, estudo de caso) e justifique a escolha.

\section{Paradigma de Pesquisa}

\section{Método de Pesquisa}

\section{Procedimentos de Coleta de Dados}

\section{Procedimentos de Análise}

\section{Aspectos Éticos}

\section{Considerações do Capítulo}

%% capitulos/5-desenvolvimento.tex

\chapter{Desenvolvimento}
\label{cap:desenvolvimento}

Descreva o artefato ou solução desenvolvida.

\section{Visão Geral da Solução}

\section{Arquitetura}

\section{Implementação}

% Exemplo de código-fonte:
\begin{lstlisting}[language=JavaScript, caption={Exemplo de código}, label={lst:exemplo}]
function hello(name) {
  console.log(`Olá, ${name}!`);
}
hello('PPGCC');
\end{lstlisting}

\section{Considerações do Capítulo}

%% capitulos/6-avaliacao.tex

\chapter{Avaliação}
\label{cap:avaliacao}

Descreva o processo de avaliação da solução desenvolvida.

\section{Planejamento da Avaliação}

\section{Instrumento de Avaliação}

\section{Participantes}

\section{Resultados}

\section{Discussão}

\section{Considerações do Capítulo}

%% capitulos/7-conclusao.tex

\chapter{Conclusão}
\label{cap:conclusao}

\section{Síntese do Trabalho}

Retome o problema de pesquisa e os objetivos, descrevendo como
foram alcançados.

\section{Contribuições}

\section{Limitações}

\section{Trabalhos Futuros}


%% ────────────────────────────────────────────────────────────
%%  ELEMENTOS PÓS-TEXTUAIS
%% ────────────────────────────────────────────────────────────
\postextual

%% Referências (obrigatório)
%% Estilo ABNT com bibtex:
\bibliographystyle{abntex2-alf}
\bibliography{pos-textuais/referencias}

%% Glossário (opcional) — descomente se necessário
% \glossary

%% Apêndices (opcional — material do próprio autor)
\begin{apendicesenv}
  \partapendices
  %% pos-textuais/apendices/apendice-a.tex
%% Apêndice: material elaborado pelo próprio autor
%% (ex.: questionários, scripts, formulários criados para a pesquisa)

\chapter{Título do Apêndice A}
\label{apend:a}

Conteúdo do apêndice elaborado pelo autor.

% Exemplo: inserir PDF de questionário
% \includepdf[pages=-]{pos-textuais/apendices/questionario.pdf}

  %\input{pos-textuais/apendices/apendice-b}
\end{apendicesenv}

%% Anexos (opcional — material de terceiros)
\begin{anexosenv}
  \partanexos
  %% pos-textuais/anexos/anexo-a.tex
%% Anexo: material de terceiros (não elaborado pelo autor)
%% (ex.: formulários de terceiros, documentos externos)

\chapter{Título do Anexo A}
\label{anex:a}

Conteúdo do anexo (material de terceiros).

% Exemplo: inserir PDF externo
% \includepdf[pages=-]{pos-textuais/anexos/documento-externo.pdf}

\end{anexosenv}

%% Índice remissivo (opcional)
% \printindex

\end{document}
